\chapter{Mixing samples without losing their meaning}
\MARSChapterHeading{Mixing samples without losing their meaning}
\label{revision:pcm-mixing}

Adding two audio streams sounds like a direct application of integer addition.
At each sample position, read one value from each stream and add them. The
problem appears when the mathematical answer does not fit the output format.
Two valid inputs can produce an invalid output value. Choosing what to do with
that value is part of the algorithm, before it becomes an instruction choice.

This chapter assumes that you can traverse an array, call a function and read
an MSA vector as a collection of integer lanes. We will use those skills to mix
two arrays of signed 16-bit samples. The result must stay in the same sample
format, and the function must report how many sums were clipped. The accompanying
lab supplies the complete calling convention and exercises memory bounds and
in-place output. Here we develop the arithmetic and the loop that make that
contract possible.

\section{A sample is a number with a storage format}

Pulse-code modulation, or PCM, represents a signal by a sequence of sample values
at regularly spaced times. Our inputs are mono: each sample position contains
one component value. A stereo stream would need a specified arrangement of left
and right components before the same array reasoning could be applied.

Each value in this exercise is a signed two's-complement halfword, from $-32768$
through $32767$. Its two bytes are stored low byte first. For example,
\texttt{ff 7f} represents $32767$, while \texttt{00 80} represents $-32768$.
The byte sequence \texttt{ff ff} represents $-1$. Loading those last two bytes as
an unsigned integer would instead give $65535$. The bits have not changed;
the interpretation used by the arithmetic has.

A scalar implementation must reconstruct a signed sample before adding it.
Because the lab permits a sample to begin at an odd byte address, two unsigned
byte loads are sufficient: combine the high byte shifted left by eight with the
low byte, then sign-extend the resulting halfword. A signed halfword load would
have the desired numeric interpretation, but it would also impose an alignment
requirement that this interface does not give us.

For inputs $a_i$ and $b_i$, let $s_i=a_i+b_i$ be the exact sum at position $i$.
Its possible range is $[-65536,65534]$. A signed 32-bit word can retain that whole
range. A signed halfword cannot. This width calculation tells us where information
could be lost: we may keep the intermediate in a word, but we must decide how to
narrow it when writing the output sample.

\section{Clipping and wrapping give different signals}

The required operation is \emph{saturation}: keep an in-range sum and replace an
out-of-range sum with the nearest representable endpoint. Thus
\[
 c_i=\min(32767,\max(-32768,s_i)).
\]
If the exact sum is $40000$, the stored sample is $32767$. If it is $-40000$,
the stored sample is $-32768$. These values are clipped. An exact sum equal to
either endpoint already fits and is not counted as clipped.

Consider the following pairs. Work out the exact sum before interpreting a
stored halfword.
\begin{center}
\tagpdfsetup{table/header-rows={1}}
\begin{tabular}{rrrrl}
\toprule
Input A & Input B & Exact sum & Output & Clipped?\\
\midrule
30000 & 2767 & 32767 & 32767 & No\\
30000 & 10000 & 40000 & 32767 & Yes\\
-30000 & -2768 & -32768 & -32768 & No\\
-30000 & -10000 & -40000 & -32768 & Yes\\
32767 & -32768 & -1 & -1 & No\\
\bottomrule
\end{tabular}
\end{center}

A store that keeps only the low sixteen bits would implement wrapping instead.
For $40000$, those bits denote $-25536$ as a signed halfword. For $-40000$,
they denote $25536$. A slightly excessive positive peak can therefore become a
large negative sample. Saturation keeps the peak at the positive endpoint.
Both operations are precisely defined, but only one meets this mixer's contract.

Clipping still changes the signal. It flattens excessive peaks; it cannot keep
the original amplitude in a format that cannot represent it. Reducing input gains
before adding is one way to avoid some clipping, but that would be a different
function. Reducing a clipped result afterward cannot recover the discarded values.

The count is another reason to retain evidence of the arithmetic. Counting output
samples equal to $32767$ would misclassify the first row of the table. To distinguish
that row from the second, we need to know whether the exact sum exceeded the bound.

\section{Use two vector results to find the clipped lanes}

An MSA vector contains 128 bits. Interpreted as halfwords, it contains eight lanes,
so one operation can mix eight independent pairs. In our pinned little-endian
MARS environment, loading sixteen consecutive sample bytes places each low/high
byte pair in the corresponding halfword lane. Changing the operation suffix
changes the lane interpretation of those bits; it does not rearrange the samples.

\texttt{adds\_s.h} adds signed halfwords with saturation. This produces the output
values we need, but does not return a count of clipped lanes. We can derive the
count by also computing \texttt{addv.h}, which keeps each sum modulo $2^{16}$.
An in-range sum has the same halfword representation in both results. An excessive
positive sum wraps into the negative half of the signed range while saturating to
$32767$. An excessive negative sum wraps into the nonnegative half while saturating
to $-32768$. For these two-input sums, the results differ exactly where clipping
occurred.

Use \texttt{ceq.h} to compare them lane by lane. An equal pair produces sixteen
one bits; an unequal pair produces sixteen zero bits. Inverting this mask makes
clipped lanes all ones. The mask is useful evidence, but an all-one halfword is
$65535$ as an unsigned number, not the integer one we want to add to a count.

Here is the complete group operation. On entry, \texttt{\$a0} and \texttt{\$a1}
point to eight readable input samples, \texttt{\$a2} points to eight writable
output samples, and \texttt{\$v1} holds the count from earlier groups. This
fragment consumes one group; pointer updates and the remaining-count test belong
to the surrounding loop. The temporary and vector registers shown are available
to the fragment.

\MARSCodeLabel{Mix eight samples and add their clipping count}
\begin{minted}{mars-idea}
ld.b $w0, 0($a0)
ld.b $w1, 0($a1)
addv.h $w2, $w0, $w1
adds_s.h $w3, $w0, $w1
ceq.h $w4, $w2, $w3
nor.v $w4, $w4, $w4
pcnt.d $w5, $w4
copy_s.d $t0, $w5[0]
copy_s.d $t1, $w5[1]
daddu $t0, $t0, $t1
dsrl $t0, $t0, 4
daddu $v1, $v1, $t0
st.b $w3, 0($a2)
\end{minted}

\texttt{pcnt.d} counts set bits separately in the two 64-bit elements of the
mask. \texttt{copy\_s.d} copies each count into a general register. There are
sixteen set bits for every clipped halfword, so adding the counts and shifting
right by four divides by sixteen. This yields a group count from zero through
eight. The shift is exact because the mask contains only whole all-zero or
all-one halfwords.

For a group whose exact sums are $0$, $32767$, $32768$, $-32768$, $-32769$,
$40000$, $-40000$ and $-1$, lanes 2, 4, 5 and 6 clip. They contribute 64 set
bits in total, so the count increases by four. The two endpoint sums contribute
nothing. Trace those lanes before using the fragment in a loop; they distinguish
the required test from a test for endpoint values.

\section{The loop must prove each access fits}

Let $r$ be the number of samples remaining. A sixteen-byte input load is justified
when $r\geq8$, because each sample occupies two bytes. After a group, advance
each pointer by sixteen and reduce $r$ by eight. The loop invariant is that all
samples before the pointers have been mixed exactly once, \texttt{\$v1} counts
only those samples, and the pointers identify the first unprocessed samples.

When fewer than eight remain, a complete vector access would extend beyond the
logical input slice. Extra backing storage does not make that access part of
the interface. Finish the tail using scalar sample loads, an exact word sum,
bounds comparisons and two byte stores. Each scalar iteration advances by two
bytes and consumes one sample. A zero-length call performs no sample access.

The same loop works when output exactly replaces either input. The group fragment
loads both inputs before storing the output. At an exact alias, that store can
replace only the current group's already loaded samples; future groups begin
after them. The scalar tail needs the same ordering for each pair. A destination
shifted into the middle of an input would need a different argument about access
order, so the lab excludes that overlap.

\begin{MARSWorkedExample}{Nine samples with in-place output}
Suppose the output pointer equals the second input pointer and the count is nine.
The first iteration reads both sixteen-byte groups before overwriting the second
input's first eight samples. One sample remains. The tail reads both final samples
before writing their result at byte offset sixteen. It must not execute another
vector load merely because the first group was aligned or the allocation is large.
\end{MARSWorkedExample}

\section{Investigate the first difference}

Compare the scalar and vector versions on identical inputs. A wrong sample and
a wrong clipping count suggest different places to investigate. If both versions
store $32767$ for an exact sum of $32767$ but only one increments its count, inspect
the clipping test. If negative inputs act like large positive numbers in the scalar
path, inspect sign extension before the sum. If the first eight samples are correct
and the ninth is wrong, inspect the tail's pointer and remaining count.

For the vector path, pause after the saturated addition, then after the comparison.
Keep the lane width explicit in the debugger. The equality mask's true halfwords
are \texttt{ffff}; interpreting the same bits as bytes or doublewords changes the
displayed grouping. After inversion and population count, expect small integer
counts in two doubleword lanes. This sequence locates the first representation
that differs from your paper trace.

Measure a complete call after checking its results and memory bounds. Vector work
has group and reduction overhead, and a short call may execute no full group at all.
A count of 63 gives seven vector groups and seven scalar iterations; a count of 64
gives eight groups and no scalar tail. Their costs need not rise smoothly with the
sample count. Native instruction counts describe guest work. Host elapsed time also
includes the simulator's costs and should be measured separately.

\begin{MARSExercise}{Change the mask reduction}
Implement the same mixer with a different method of counting clipped lanes. Keep
the sample results, zero-length behavior, memory bounds and permitted aliases.
Explain how your method converts an all-one comparison lane into one contribution
to the count. Test exact endpoints as well as values on both sides of each bound.
Use lengths 0, 1, 7, 8, 9 and 16 to expose group and tail mistakes.
\end{MARSExercise}

\begin{MARSExercise}{Decide when to clip}
Mix the three values $30000$, $10000$ and $-10000$. Compare clipping after each
pairwise addition with retaining an exact wider total and clipping once at the end.
Explain which information the first clipping step discards. State an accumulation
policy before extending the two-input mixer to more streams.
\end{MARSExercise}

\section*{Sources and laboratory connection}
The numerical representation follows the signed 16-bit PCM format described in
Microsoft's \href{https://learn.microsoft.com/en-us/windows/win32/multimedia/devices-and-data-types}{Devices and data types}.
The particular clipping count, arbitrary alignment and allowed aliases are the
contract of the PCM mixing lab. The MSA operations above are checked against the
pinned MARS integer and register instruction implementations. The public
\path{packages/020-pcm-mixing/} contains the caller, input recipes and listening
comparisons. The written lab supplement develops the observations needed to build
and check the complete function; this fragment alone is not that function.
